\section{Smooth profiles with a pressurized core}\label{sec:profiles}

We work in dimension $12$.  The small parameter is $a>0$; later we fix
$a=(p-1)^{-1}$ with $p$ an odd integer.  In this section $a$ may take every
sufficiently small positive real value.  The profiles never vanish, so the
nonlinearity is smooth along them even without the restriction on $p$.
We use the parameter disk from Proposition~\ref{free:main}; all
terminating decimals below denote exact rational numbers.  Thus
\[
 \mathcal D=\{(b,Z):|(b,Z)-(b_*,Z_*)|\le10^{-8}\},\qquad
 b_*=.33477606871236,\quad Z_*=2.70506819293654.
\]
Put
\begin{equation}\label{profile:radii}
 \varepsilon=10^{-4},\qquad r_b=2\sqrt{Z_*}+\varepsilon,
 \qquad R(Z)=2\sqrt Z,\qquad r_0=r_b-6\varepsilon.
\end{equation}
The radius $r_b$ is fixed throughout the construction.

\subsection{Similarity variables and the amplitude equations}

For $\tau=T-t$, $s=\log(T/\tau)$ and
$y=(x-x_*)/\sqrt\tau$, write
\[
 u(t,x)=e^{i\theta}\tau^{-a}(\tau/T)^{ib}U(s,y).
\]
Since $\partial_t s=\tau^{-1}$ and
$\partial_t y=y/(2\tau)$, substitution into the defocusing equation gives
\begin{equation}\label{profile:similarity}
 i\partial_sU+\Delta U+i\left(\frac{y\cdot\nabla}{2}+a\right)U+bU
       =|U|^{1/a}U.
\end{equation}
We seek a nonvanishing radial stationary solution $Q=Ae^{i\phi}$, with
$A>0$ and $\phi(0)=0$.  Throughout this section a prime denotes a radial
derivative, except when explicitly attached to the last argument of $H$.
Set
\begin{equation}\label{profile:notation}
 P=A^{1/a},\quad \mu=A^2,\quad c_*=6-2a,\quad
 w=\frac r2+2\phi',\quad D=r\partial_r.
\end{equation}
The same letter $w$ denotes the vector field $w(r)y/r$, extended smoothly
at zero, and $W=w\cdot\nabla$.  The imaginary part of the stationary
equation is
\[
 2\frac{A'}A\phi'+\phi''+\frac{11}{r}\phi'
       +\frac r2\frac{A'}A+a=0.
\]
It follows that
\begin{equation}\label{profile:amplitude}
 \begin{split}
 w(r)&=c_*r^{-11}A(r)^{-2}\int_0^r A(t)^2t^{11}\,dt,\\
 -\Delta A+A^{1+1/a}&=VA,\qquad
 V=b+\frac{r^2}{16}-\frac{w^2}{4},\\
 w'&=c_*-\frac{11w}{r}-2w\frac{A'}A.
 \end{split}
\end{equation}
Here the integration constant in the first line is zero by regularity at
the origin.  The real part gives the second line because
$\phi'^2+r\phi'/2=w^2/4-r^2/16$.  In particular,
\begin{equation}\label{profile:divergence}
 \operatorname{div}(\mu w)=c_*\mu.
\end{equation}
Conversely, positive amplitudes satisfying the first two equations in
\eqref{profile:amplitude}, with $A'(0)=0$, give a stationary solution by
integrating $\phi'=(w-r/2)/2$.

Here is the complete profile statement used subsequently.  A constant
with a subscript $j$ may depend on that fixed derivative order.  Uniform
statements as $a\downarrow0$ refer to all profiles selected by the
construction, without requiring a continuous selection of their matching
parameters.

\begin{proposition}[Profiles and uniform estimates]\label{profile:existence}
There is $a_0>0$ such that, for every $0<a<a_0$, there are
$(b_a,Z_a)\in\mathcal D$ and a smooth nonvanishing radial stationary
solution $Q_a=A_ae^{i\phi_a}$ of \eqref{profile:similarity}, normalized by
$Q_a(0)>0$.  They have the following properties.
\begin{enumerate}
\item The identities \eqref{profile:amplitude}--\eqref{profile:divergence}
hold, $0<P_a\le1$, and there are constants $c,C>0$, independent of small
$a$, such that
\begin{equation}\label{profile:coercive-bounds}
 C^{-1}\langle r\rangle^{-4a}\le\mu_a(r)
       \le C\langle r\rangle^{-4a},\qquad
 \nabla w_a\ge c I.
\end{equation}
\item If $a_n\downarrow0$ and $(b_{a_n},Z_{a_n})\to(b,Z)$, then
$j(b,Z)=0$, where $j$ is the free matching function of
Proposition~\ref{free:main}.  With $R=2\sqrt Z$, the profiles converge in
$C^1$ on every compact ball to
\begin{equation}\label{profile:limit}
 Q_0(r)=
 \begin{cases}
 1,&0\le r\le R,\\[2pt]
 \displaystyle\frac{H(-ib,6,-ir^2/4)}{H(-ib,6,-iZ)},&r\ge R.
 \end{cases}
\end{equation}
The exterior amplitude is strictly decreasing and has a positive limit
at infinity.  The convergence is $C^\infty$ on compact subsets of
$\{r>R\}$.  On every fixed ball,
\begin{equation}\label{profile:pressure-limit}
 P_{a_n}\ \stackrel{*}{\rightharpoonup}\ b\mathbf1_{\{r<R\}}
       \quad\hbox{in }L^\infty,
\end{equation}
and $\mu_{a_n},w_{a_n}$ converge in $C^1$.  In particular
$P_{a_n}/a_n\to0$ locally uniformly on $\{r>R\}$.
\item On the fixed exterior $r\ge r_b$,
\begin{equation}\label{profile:tail}
 Q_a(r)=r^{\nu_a}f_a(r),\qquad \nu_a=-2a+2ib_a,
 \qquad 0<\kappa\le|f_a(r)|\le\rho<1,
\end{equation}
with $\kappa,\rho$ independent of $a$.  There are coefficients
$m_{a,j}$, with $|m_{a,0}|\ge\kappa$, such that for every pair of fixed
nonnegative integers $N,J$,
\begin{equation}\label{profile:symbol-expansion}
 \left|D^j\left(f_a(r)-\sum_{q=0}^{N}m_{a,q}r^{-2q}\right)\right|
       \le C_{N,J}r^{-2N-2},\qquad 0\le j\le J,\quad r\ge r_b.
\end{equation}
The constants are uniform as $a\downarrow0$.  In particular,
\begin{equation}\label{profile:exterior-pressure}
 |D^j(P_a/a)|\le\eta_j(a)r^{-2},\qquad
 |D^je^{2i\phi_a}|\le C_j,\qquad \eta_j(a)\longrightarrow0
\end{equation}
for each fixed $j$.  On the inner ball,
\begin{equation}\label{profile:pressure-transport}
 W_a(P_a/a)\le Ca(P_a/a),\qquad r\le r_b,
\end{equation}
whereas $|W_a(P_a/a)|\le\eta(a)r^{-2}$ for $r\ge r_b$, with
$\eta(a)\to0$.  Consequently the positive part of $W_a(P_a/a)$ is
bounded uniformly on the whole space.
\item For each fixed $a$, every Cartesian derivative obeys
\begin{equation}\label{profile:fixed-power-symbols}
 |\partial^\beta Q_a(y)|\le C_{a,\beta}\langle y\rangle^{-2a-|\beta|},
 \qquad
 \left|\partial^\beta\big((P_a/a)e^{2i\phi_a}\big)\right|
       +|\partial^\beta(P_a/a)|
       \le C_{a,\beta}\langle y\rangle^{-2-|\beta|}.
\end{equation}
The time-symmetry combination has the improved exterior order
\begin{equation}\label{profile:time-symbol}
 \left|D^j\left((a-ib_a+D/2)Q_a\right)\right|
       \le C_j r^{-2a-2},\qquad r\ge r_b.
\end{equation}
\end{enumerate}
\end{proposition}

We prove the proposition in stages.  The inner problem is solved on the
whole parameter disk, as is the slow exterior problem.  Only their final
matching uses a choice of $(b_a,Z_a)$.

\subsection{The inner boundary value problem}

Let $F(r;b,R')$ solve
\begin{equation}\label{profile:free-initial}
 F''+\left(\frac{11}{r}+\frac{ir}{2}\right)F'+bF=0,
       \qquad F(R')=1,\quad F'(R')=0,
\end{equation}
and set
\[
 F_{\rm in}(r;b,Z)=F(r;b,R(Z)),\qquad
 h_b(b,Z)=|F_{\rm in}(r_b;b,Z)|.
\]
The subscript on $h_b$ indicates the fixed boundary radius, not
dependence on $b$ alone.

\begin{lemma}[The boundary modulus]\label{profile:boundary-modulus}
Uniformly on $\mathcal D$,
\begin{equation}\label{profile:boundary-gap}
 1-.2\varepsilon^2<h_b<1-.15\varepsilon^2.
\end{equation}
For fixed $b$ in the projection of $\mathcal D$, the function
$R'\mapsto|F(r_b;b,R')|$ is strictly increasing on $[r_0,r_b)$.
\end{lemma}

\begin{proof}
On these intervals the coefficient $11/r+ir/2$ has modulus less than
$6$, and $.334<b<.335$.  Put $u=r-R'$, so $0\le u\le6\varepsilon$.
Integrating the two first-order equations in \eqref{profile:free-initial}
gives, by a bootstrap starting at $u=0$,
\[
 |F'|\le .35u,\qquad |F-1|\le .175u^2,
       \qquad |F''+b|<.003.
\]
For completeness, assuming the first bound up to a first failure gives
$|F-1|\le.175u^2$ and
$|F''+b|\le6(.35)u+.335(.175)u^2<.0013$; integration then gives
$|F'|<.337u<.35u$, excluding that failure.  The displayed weaker
$.003$ bound follows as well.  Thus
\[
 F(r_b)=1-\frac b2(r_b-R')^2+E,\qquad
       |E|\le.0015(r_b-R')^2.
\]
Since $|R(Z)-2\sqrt{Z_*}|<10^{-8}$, taking $R'=R(Z)$ proves
\eqref{profile:boundary-gap}, including the negligible quadratic error
when passing from the complex value to its modulus.

Differentiate the initial-value problem with respect to $R'$.  The
derivative $G=\partial_{R'}F$ satisfies the same linear equation and has
data $G(R')=0$, $G'(R')=b$.  Integration, or the same bootstrap, gives
$|G(r_b)/(r_b-R')-b|<.003$.  As $F(r_b)=1+O((r_b-R')^2)$,
\[
 \partial_{R'}|F(r_b)|
   =\frac{\operatorname{Re}(\overline{F(r_b)}G(r_b))}{|F(r_b)|}>0.
\]
\end{proof}

\begin{lemma}[The inner amplitude]\label{profile:inner}
For all sufficiently small $a>0$ and all $(b,Z)\in\mathcal D$, the
system \eqref{profile:amplitude} has a positive radial solution on
$[0,r_b]$ with $A'(0)=0$ and $A(r_b)=h_b$.  Within the bounds
\begin{equation}\label{profile:inner-bounds}
 h_b\le A\le .5^a,
\end{equation}
it is unique.  It depends continuously on $(b,Z)$, and obeys
\begin{equation}\label{profile:bulk-barrier}
 A(r)\ge .2^a-.02(r-r_0)_+^2.
\end{equation}
The amplitudes are bounded uniformly in $C^{1,1}[0,r_b]$.
\end{lemma}

\begin{proof}
First take $a$ small enough that $h_b^{1/a}\le.3$ uniformly on the disk;
this also ensures $h_b<.5^a$.  For a continuous radial input $A$ satisfying
\eqref{profile:inner-bounds}, define $w,V$ by the first and second lines
of \eqref{profile:amplitude}.  The useful nonsingular formula is
\begin{equation}\label{profile:w-ratio}
 \frac{w(r)}r=c_* A(r)^{-2}\int_0^1A(rs)^2s^{11}\,ds.
\end{equation}
Since $h_b>.999$, $r_b^2<11$, and $a$ is small, it implies
\begin{equation}\label{profile:V-bounds}
 .3\le V\le .5,\qquad
 \|V[A_1]-V[A_2]\|_\infty\le10\|A_1-A_2\|_\infty.
\end{equation}
Indeed $c_*h_b^2/12\le w/r\le c_*/(12h_b^2)$; inserting these bounds
into $V=b+r^2(1/16-(w/r)^2/4)$ gives the first assertion for small $a$.
The Lipschitz constant of $A\mapsto w/r$ is at most
$c_*(2h_b^{-3}+2h_b^{-2})/12<2.01$; multiplying by
$r_b^2(\sup w/r)/2$ gives a constant less than $10$ for $V$.

For this prescribed $V$, solve the scalar Dirichlet problem
\begin{equation}\label{profile:scalar-dirichlet}
 -\Delta H+H^{1+1/a}=VH,\qquad H(r_b)=h_b,\quad H'(0)=0.
\end{equation}
The constants $h_b$ and $.5^a$ are respectively a lower and an upper
solution.  One can obtain $H$ by monotone iteration: choose $L>0$ so that
$t\mapsto(L+V(r))t-t^{1+1/a}$ is increasing on $[h_b,.5^a]$, and at
each step invert $-\Delta+L$ with the given boundary value.  This inverse
is positive by the minimum principle.  Its existence, if needed directly
in the radial class, follows by minimizing the strictly convex functional
$\int(|\nabla u|^2+L|u|^2-2fu)/2$ for zero boundary data; the radial
equation then reads $(r^{11}u')'=r^{11}(Lu-f)$ and gives the classical
solution.  Iteration from $h_b$ is increasing and stays below $.5^a$.
For fixed $a$ its right sides are uniformly bounded; radial integration
gives compactness in $C^1$, and its limit solves
\eqref{profile:scalar-dirichlet}.

We record a comparison argument that also applies below.  The radial
Dirichlet inverse of $-\Delta$ is
\[
 (Kf)(r)=\int_r^{r_b}t^{-11}\int_0^t s^{11}f(s)\,ds\,dt,
       \qquad \|K\|_{\infty\to\infty}=r_b^2/24<11/24.
\]
Thus the positive geometric series defines
$T=(-\Delta-.5)^{-1}=\sum_{n\ge0}(.5K)^nK$ with zero boundary values.
Also $\psi=1+K(1+.5\psi)$ has a positive solution and satisfies
$(-\Delta-.5)\psi=1$, $\psi(r_b)=1$.  For every $q\ge0$ and $V\le.5$,
\[
 (-\Delta-V+q)\psi\ge1.
\]
If a function with nonnegative boundary values violated the comparison
principle for $-\Delta-V+q$, its ratio to $\psi$ would have a negative
interior minimum, where the product rule gives a strictly negative value
for that operator.  This is a contradiction.  Taking $q$ to be the
divided difference of $t^{1+1/a}$ proves uniqueness in
\eqref{profile:scalar-dirichlet} and comparison with lower solutions.

The function
$B(r)=.2^a-.02(r-r_0)_+^2$ is such a lower solution.  For small $a$ it
is positive, $B(r_b)<h_b$ by \eqref{profile:boundary-gap},
$B^{1/a}\le.2$, and
\[
 -\Delta B\le .04+.44\frac{6\varepsilon}{r_0}<.05
       \le(.3-.2)B.
\]
At $r_0$ the first derivative is continuous, so no interface measure
occurs in this inequality.  Comparison proves
\eqref{profile:bulk-barrier} for every output $H$.

Let $H_1,H_2$ be outputs for two inputs.  Their difference solves
\[
 (-\Delta-V_1+q)(H_1-H_2)=(V_1-V_2)H_2,
 \quad
 q=\int_0^1(1+1/a)(tH_1+(1-t)H_2)^{1/a}\,dt.
\]
The barrier implies $q\ge.2/a$ on $r\le r_0$, whereas $q\ge0$
everywhere.  A positive supersolution for unit forcing is
\begin{equation}\label{profile:small-inverse}
 S=10a+2T\mathbf1_{\{r>r_0\}}.
\end{equation}
On the inner ball, the constant term contributes at least $2-5a$ and
the second term contributes a nonnegative quantity.  On the shell, the
second term contributes at least $2$, while the constant term is bounded
below by $-5a$.  Thus the comparison principle applies, including at
the shell interface.  To bound $S$, radial integration gives
\[
 \|K\mathbf1_{\{r>r_0\}}\|_\infty
 \le\int_{r_0}^{r_b}(t-r_0)\,dt=(6\varepsilon)^2/2,
 \qquad
 \|T\mathbf1_{\{r>r_0\}}\|_\infty\le(6\varepsilon)^2.
\]
Here the second inequality uses
$(1-.5\|K\|)^{-1}<2$.  Consequently the input-output map satisfies
\[
 \|H_1-H_2\|_\infty
 \le10\big(10a+2(6\varepsilon)^2\big)\|A_1-A_2\|_\infty.
\]
Our already fixed shell width makes the second contribution less than
$1/4$; taking $a$ small makes the whole factor less than $1/2$.
The closed interval \eqref{profile:inner-bounds} in the radial continuous
functions is complete, so successive iteration gives its unique fixed
point.  This proves the coupled inner problem.

Finally $0<P\le.5$ and \eqref{profile:V-bounds} give
\[
 A'(r)=r^{-11}\int_0^r t^{11}(P(t)-V(t))A(t)\,dt,
       \qquad |A'|\le Cr,\quad |A''|\le C.
\]
These are the asserted uniform bounds.  For fixed $a$, convergent
parameter sequences have $C^1$ convergent subsequences; the equation and
uniqueness identify their limits.  This proves continuous dependence,
including when the boundary value $h_b$ varies.
\end{proof}

\begin{lemma}[Identification of the core]\label{profile:core-limit}
For the inner solutions of Lemma~\ref{profile:inner}, suppose
$a_n\downarrow0$ and $(b_n,Z_n)\to(b,Z)\in\mathcal D$.  Then their
complex reconstructions converge in $C^1[0,r_b]$ to the function equal to
$1$ on $[0,R(Z)]$ and to $F_{\rm in}$ on $[R(Z),r_b]$.
Their pressures converge weak star to $b\mathbf1_{\{r<R(Z)\}}$.
These convergence assertions are uniform over the compact parameter disk
in the sequential sense just stated.
\end{lemma}

\begin{proof}
The preceding bounds give $C^1$ compactness of the amplitudes.  Every
limit $A_0$ is at most $1$ and, by \eqref{profile:bulk-barrier}, equals
$1$ on $[0,r_0]$.  Formula \eqref{profile:w-ratio} and the last equation
in \eqref{profile:amplitude} give $C^1$ convergence of $w$, including at
zero.  On every compact interval where $A_0<1$, the pressure tends
uniformly to zero.  There the limit satisfies $\Delta A_0=-V_0A_0$,
with $V_0\ge.3$.

Each component of $\{A_0<1\}$ has a left endpoint $R'$ where
$A_0(R')=1$, $A_0'(R')=0$.  On this component
$(r^{11}A_0')'=-r^{11}V_0A_0<0$, so $A_0'<0$ and the component cannot
have a right endpoint below $r_b$.  As $A_0(r_b)=h_b<1$, there is exactly
one such component, $(R',r_b]$, and $R'\in[r_0,r_b)$.  On the plateau
\eqref{profile:w-ratio} gives $w_0=r/2$ and hence $\phi_0=0$.
Reconstruction gives the free complex equation with data $1,0$ at $R'$.
Its boundary modulus is $h_b$, so
Lemma~\ref{profile:boundary-modulus} forces $R'=R(Z)$.

The pressures are bounded by $.5$.  Any weak-star limit $P_0$ satisfies
in distributions
\[
 P_0A_0=V_0A_0+\Delta A_0.
\]
There is no surface measure at $R$ because $A_0$ and $A_0'$ are
continuous.  On the plateau the right side is $b$, and on its exterior
it is zero.  Since $A_0$ is positive, this identifies $P_0$ uniquely.
The preceding identification applies to every subsequence; compactness
therefore proves all convergence assertions.  Reconstruction by
$\phi'=(w-r/2)/2$ gives the asserted $C^1$ convergence of the complex
profiles as well.
\end{proof}

\begin{figure}[tbp]
  \centering
  \begingroup
  \definecolor{coreblue}{RGB}{48,91,120}
  \definecolor{coreorange}{RGB}{163,92,43}
  \resizebox{0.95\linewidth}{!}{%
  \begin{tikzpicture}[
    x=1cm,y=1cm,font=\small,
    axis/.style={->,thin,black!75},
    guide/.style={densely dashed,thin,black!45},
    amplitude/.style={very thick,coreblue},
    pressure/.style={very thick,coreorange},
    every node/.style={inner sep=2pt}]
    \begin{scope}
      \fill[coreorange!7] (0,0) rectangle (2.2,2.6);
      \draw[axis] (0,0) -- (5.9,0) node[right] {$r$};
      \draw[axis] (0,0) -- (0,2.8) node[above] {$A_0(r)$};
      \draw[guide] (2.2,0) -- (2.2,2.55);
      \draw[amplitude] (0,2.1) -- (2.2,2.1)
        .. controls (2.85,2.1) and (3.15,1.3) .. (3.85,1.02)
        .. controls (4.4,.8) and (5.0,.7) .. (5.65,.66);
      \draw (0,2.1) -- (-.08,2.1) node[left] {$1$};
      \draw (2.2,0) -- (2.2,-.09) node[below] {$R$};
      \node[below left] at (0,0) {$0$};
      \node at (1.1,1.15) {constant core};
      \node[align=center] at (4.3,2.25) {free exterior\\strictly decreasing};
      \node at (2.85,-.65) {(a) Limiting amplitude};
    \end{scope}
    \begin{scope}[xshift=7.4cm]
      \fill[coreorange!7] (0,0) rectangle (2.2,2.6);
      \draw[axis] (0,0) -- (5.9,0) node[right] {$r$};
      \draw[axis] (0,0) -- (0,2.8) node[above] {$P_0(r)$};
      \draw[guide] (2.2,0) -- (2.2,2.55);
      \draw[pressure] (0,2.1) -- (2.2,2.1);
      \draw[pressure] (2.2,0) -- (5.65,0);
      \draw[pressure,fill=white] (2.2,2.1) circle (.045);
      \draw[pressure,fill=white] (2.2,0) circle (.045);
      \draw (0,2.1) -- (-.08,2.1) node[left] {$b$};
      \draw (2.2,0) -- (2.2,-.09) node[below] {$R$};
      \node[below left] at (0,0) {$0$};
      \node at (1.1,1.15) {pressure $b$};
      \node at (4.15,1.15) {zero pressure};
      \node at (2.85,-.65) {(b) Limiting pressure};
    \end{scope}
  \end{tikzpicture}%
  }
  \endgroup
  \caption{Schematic of the singular profile limit as $a\to0$.
  The amplitude $A_0$ equals $1$ for $r\le R$ and decreases strictly
  on the free exterior; its limiting tail is nonzero. The pressure
  converges weakly to $P_0=b\mathbf 1_{\{r<R\}}$; its value at the
  jump is immaterial. Compact support of this limiting pressure does
  not assert compact support of the profiles $Q$.
  The spectral interface conditions for amplitude and phase perturbations
  are derived in Lemma~\ref{spec:interface}.
  Heights and radial distances are schematic.}
  \label{fig:core-limit}
\end{figure}


The pressure limit in Figure~\ref{fig:core-limit} is compatible with
$A_a\to1$ in the core: the power $1/a$ retains the order-$a$ variation of
the amplitude.  Weak-star convergence, rather than differentiation of
this limiting step function, is the fact needed in the constrained mode
problem.

For each $(b,Z)\in\mathcal D$, the inner solution now has prescribed
boundary modulus $h_b$ and a determined complex logarithmic derivative
at $r_b$. We next construct a slow exterior solution with the same
boundary modulus. A constant phase rotation can then match the boundary
values; equality of the logarithmic derivatives supplies the two
remaining real matching conditions. Lemma~\ref{profile:matching} will
impose these conditions by varying $(b,Z)$.

\subsection{A uniform slow solution on the exterior}

Write
\[
 G_b(r)=H(-ib,6,-ir^2/4),\qquad
 j(b,z)=-\frac{\partial_xH(-ib,6,-iz)}{H(-ib,6,-iz)}.
\]
Proposition~\ref{free:main} says that $G_b$ never vanishes for
$r\ge R(Z)$.  For its amplitude and phase,
\begin{equation}\label{profile:free-j}
 \frac{A'}A=-\frac r2\operatorname{Im}j(b,r^2/4),\qquad
 \frac wr=\frac12+\operatorname{Re}j(b,r^2/4)<.543.
\end{equation}
At $R$, $w/r>0$ because $|j(b,Z)|\le3\cdot10^{-9}$.
The divergence equation with $c_*=6$ shows that $w$ stays positive.
For $R\le r\le2\sqrt3$, it follows that
\[
 V=b+r^2/16-w^2/4>.18.
\]
Normalize temporarily by $|G_b(R)|=1$.  Its initial amplitude derivative
has magnitude at most $(R/2)3\cdot10^{-9}$, and the equation
$(r^{11}A')'=-r^{11}VA$ makes $A'(r_b)<0$: on $[R,r_b]$ the amplitude
stays, for example, between $1/2$ and $2$, while integrating the negative
right side decreases $r^{11}A'$ by more than its initial magnitude.
It continues to decrease up to $2\sqrt3$.  Beyond that radius the strict
inequality $\operatorname{Im}j>0$ proves $A'<0$ directly.  Thus $|G_b|$
is strictly decreasing on $[r_b,\infty)$, uniformly over the disk in the
sense of strict negativity on each compact interval.  If $j(b,Z)=0$,
the same argument gives strict decrease throughout $(R,\infty)$.

For reference define the normalized free exterior
\begin{equation}\label{profile:normalized-free}
 Q^{\rm out}_0(r;b,Z)=h_b(b,Z)\frac{G_b(r)}{G_b(r_b)}.
\end{equation}
It has positive real value $h_b$ at $r_b$.  The free symbol expansion
shows that $r^{-2ib}Q^{\rm out}_0$ tends to a nonzero number $m_0(b,Z)$.
By the preceding monotonicity its modulus is at most $h_b<1$ everywhere.
Nonvanishing, the nonzero limit, and compactness of the parameter disk
therefore give uniform lower and upper bounds strictly between $0$ and
$1$ for its modulus.

\begin{lemma}[Slow exterior construction]\label{profile:exterior}
For all sufficiently small $a>0$ and $(b,Z)\in\mathcal D$, there is a
nonvanishing exterior stationary solution $Q^{\rm out}_a$ with
$Q^{\rm out}_a(r_b)=h_b$.  It depends continuously on $(a,b,Z)$, down to
the limit \eqref{profile:normalized-free}, with convergence in every
fixed derivative on compact exterior intervals.  Its factor
$f=r^{2a-2ib}Q^{\rm out}_a$ satisfies the bounds and expansions in
\eqref{profile:tail}--\eqref{profile:symbol-expansion}, uniformly over
the entire parameter disk.
\end{lemma}

\begin{proof}
Writing $\nu=-2a+2ib$ and $Q=r^\nu f$ transforms the equation into
\begin{equation}\label{profile:f-equation}
 (D+\nu)(D+\nu+10)f+\frac{ir^2}{2}Df=N_a(f),
       \qquad N_a(f)=|f|^{1/a}f.
\end{equation}
Set $N_0=0$.  Choose a small fixed neighborhood of the compact family
$(m_0(b,Z),b)$ such that the corresponding free functions $f$ stay in
an annulus
\begin{equation}\label{profile:annulus}
 0<2\kappa\le|f|\le\rho_0<1,\qquad r\ge r_b.
\end{equation}
This is possible because at $a=0$ the equation is complex linear in the
coefficient $m=\lim f$.  Enlarge the annulus slightly, still keeping its
upper radius below $1$.  On this enlarged annulus every fixed real
derivative of $N_a$ tends uniformly to zero as $a\downarrow0$: its norm
is bounded by $C_j a^{-j}\rho_1^{1/a}$ for some $\rho_1<1$ (adjusting
$C_j$ and the power of $a$ does not affect this conclusion).  Thus all
fixed differentiability orders needed below are uniformly controlled.

For prescribed $m,b$, construct a formal expansion
$f_N=\sum_{q=0}^Nm_qr^{-2q}$ with $m_0=m$.  Substituting in
\eqref{profile:f-equation} determines $m_{q+1}$ from the coefficient of
$r^{-2q}$: its coefficient in $(ir^2/2)Df_N$ is $-i(q+1)m_{q+1}$.
All the other terms at that order involve only $m_0,\ldots,m_q$ and
derivatives of $N_a$ at $m$.  This proves existence and uniqueness of
the coefficients, boundedness at every fixed order, continuous
dependence down to $a=0$, and the same bounds after a fixed number of
derivatives in the real and imaginary parts of $m$ and in $b$.
The residual of $f_N$ is $O(r^{-2N})$, with its fixed symbol and
parameter derivatives.

Here is a convergent construction behind this expansion.  Put $t=\log r$
and $Y=(e,De)$ for $e=f-f_N$.  The error equation has the form
\begin{equation}\label{profile:tail-system}
 \partial_tY=
 \begin{pmatrix}0&0\\0&-ie^{2t}/2\end{pmatrix}Y
       +\mathcal B Y+\mathcal N(t,Y)+\mathcal R_N(t),
\end{equation}
where the constant matrix $\mathcal B$ has bounded norm,
$\mathcal N(t,0)=0$, its Lipschitz constant on the annulus is bounded
uniformly in $a,m,b$, and
$|\mathcal R_N(t)|\le C_Ne^{-2Nt}$.  The propagator $\Phi(t,s)$ of the
displayed diagonal part is unitary.  The backward integral equation is
\begin{equation}\label{profile:tail-contraction}
 Y(t)=-\int_t^\infty\Phi(t,s)
       \big(\mathcal B Y(s)+\mathcal N(s,Y(s))+\mathcal R_N(s)\big)\,ds.
\end{equation}
On $t\ge t_*$ use the norm
$\|Y\|_N=\sup e^{2Nt}|Y(t)|$.  If $C$ bounds the combined Lipschitz
constant of the two terms involving $Y$, their integral has norm at
most $C/(2N)$.  Choose one integer $N_*$ with this ratio below $1/2$,
and then choose $t_*$ large enough that the corresponding ball around
$f_{N_*}$ remains in the enlarged annulus.  The residual term has norm
at most $C_{N_*}/(2N_*)$, so iteration in a ball of twice this radius
converges geometrically to a solution.  No decay estimate for an
oscillatory propagator has been used; unitarity and the weight make the
integral absolutely convergent.

The same construction works at any larger truncation order.  Two
solutions so obtained differ by $O(r^{-2N_*})$ in both entries after
subtracting their common expansion through order $N_*$.  Their
difference solves \eqref{profile:tail-contraction} without the residual,
on a common sufficiently long tail.  The contraction estimate forces
it to vanish there, and uniqueness for the finite-interval ODE makes
the solutions identical.  This proves an expansion to every order for
the single solution already constructed.  Symbol derivatives cause no
loss in this conclusion.  Indeed \eqref{profile:tail-system} bounds the
next derivative of an error by at most $C r^2$ times its two entries and
the residual; differentiating it $J$ times loses at most $2J$ powers of
$r$.  Constructing at an order greater than $N+J+2$ first therefore
gives the remainder estimate of \eqref{profile:symbol-expansion} through
the prescribed order $J$.  This also explains why the constants need
only be uniform at each fixed order.

Differentiating the integral equation with respect to $m$ or $b$ gives
the same contraction for its derivative, with a bounded inhomogeneous
term.  Difference quotients converge to that solution, proving real
$C^1$ dependence; repeated differentiation works at every fixed order.
Continuity at $a=0$ follows directly from the uniform contraction and
convergence of its coefficients.  Continue the solution from the fixed
large radius back to $r_b$.  On this finite interval the right side of
the first-order form of \eqref{profile:f-equation} and all its fixed
derivatives are bounded and converge on the enlarged annulus.  The
integral equation on finitely many short intervals, followed by
Gronwall's inequality for differences, keeps the solution close to its
free counterpart and proves existence and parameter convergence down
to $r_b$.  This also proves convergence of all fixed radial derivatives
and the first parameter derivatives there.  In view of
\eqref{profile:annulus}, the final functions satisfy
$\kappa\le|f|\le\rho<1$ for fixed $\kappa,\rho$, uniformly in small
$a$ and the chosen parameter neighborhood.

It remains to choose $m$ so that $Q(r_b)=h_b$.  At $a=0$ the map
$m\mapsto Q(r_b)$ is complex multiplication by a nonzero factor.
Its real two-dimensional inverse is uniformly bounded over the compact
parameter disk.  The just proved $C^1$ convergence implies that, on a
fixed small neighborhood of each $m_0$, the equation for $m$ is a
contraction after multiplying by that limiting inverse.  Its residual
at $m_0$ tends to zero uniformly, so it has a unique root near $m_0$.
These roots agree on overlapping neighborhoods and depend continuously
on $(a,b,Z)$.  This imposes the required boundary value.  All tail bounds
extend from the initially fixed large radius to $r_b$ by the compact
interval bounds, completing the proof.
\end{proof}

\subsection{Matching, pressure transport, and deformation}

\begin{lemma}[Matching by degree]\label{profile:matching}
For every sufficiently small $a>0$ there is $(b_a,Z_a)\in\mathcal D$
at which the inner and outer profiles have the same logarithmic
derivative at $r_b$.  After a constant phase rotation of the outer
profile they form a smooth nonvanishing stationary profile.  Every
parameter limit as $a\downarrow0$ is a zero of $j$ in $\mathcal D$.
\end{lemma}

\begin{proof}
Let $Q^{\rm in}_a$ be reconstructed with phase zero at the origin, and
define the continuous complex-valued matching map
\[
 M_a(b,Z)=\frac{(Q^{\rm in}_a)'(r_b)}{Q^{\rm in}_a(r_b)}
       -\frac{(Q^{\rm out}_a)'(r_b)}{Q^{\rm out}_a(r_b)}.
\]
Lemmas~\ref{profile:core-limit} and \ref{profile:exterior} give uniform
convergence on $\mathcal D$ to
$M_0=F_{\rm in}'(r_b)/F_{\rm in}(r_b)-G_b'(r_b)/G_b(r_b)$.
Both denominators are nonzero.

The Wronskian $\mathcal W=F_{\rm in}G_b'-F_{\rm in}'G_b$ satisfies
$\mathcal W'=-(11/r+ir/2)\mathcal W$.  At $R=2\sqrt Z$,
$G_b'(R)=(iR/2)j(b,Z)G_b(R)$, and hence
\begin{equation}\label{profile:wronskian-match}
 M_0(b,Z)=K(b,Z)j(b,Z),\qquad
 K=-\frac{iR}{2}\left(\frac R{r_b}\right)^{11}
 e^{-i(r_b^2-R^2)/4}
 \frac{G_b(R)}{F_{\rm in}(r_b)G_b(r_b)}.
\end{equation}
The factor $K$ is continuous and never zero on the entire disk.
Contracting its parameter argument to the disk center gives a homotopy
through nonzero factors to a constant.  Complex multiplication by that
constant preserves orientation.  Proposition~\ref{free:main} therefore
gives
\[
 \deg(M_0,\mathcal D^\circ,0)
       =\deg(j,\mathcal D^\circ,0)\ne0.
\]
In particular $M_0$ is separated from zero on the boundary.  Uniform
convergence preserves the boundary homotopy and the degree for every
sufficiently small $a$, so $M_a$ has a zero.

At such a zero the values have equal modulus $h_b$.  Rotate the outer
profile to match the value of the inner one; equality of logarithmic
derivatives then matches the first derivatives as well.  The ODE gives
smoothness across $r_b$, since the profile is nonzero.  To verify
smoothness through the origin without a regular-singular shortcut,
write $q(x)=Q(\sqrt x)$ for $x\ge0$.  Integrating the complex radial
equation gives
\[
 q'(x)=\frac14\int_0^1t^5e^{-ix(1-t)/4}
       \big(|q(xt)|^{1/a}-b-ia\big)q(xt)\,dt.
\]
The right side is continuous at $x=0$.  It first proves $q\in C^1$ and
then, by successive differentiation, $q\in C^\infty$ because $q$ is
bounded away from zero on this compact interval.  Thus
$Q(y)=q(|y|^2)$ is smooth at the origin.  The normalization at zero was
already imposed.  Finally, a convergent sequence of matching parameters
satisfies $M_0(b,Z)=0$, so \eqref{profile:wronskian-match} gives
$j(b,Z)=0$.  At such a root the free initial-value solution equals
$G_b/G_b(R)$, proving \eqref{profile:limit}.
\end{proof}

\begin{lemma}[The pressure derivative and positive deformation]
\label{profile:pressure}
The matched profiles satisfy \eqref{profile:coercive-bounds} and
\eqref{profile:exterior-pressure}--\eqref{profile:pressure-transport}.
\end{lemma}

\begin{proof}
First consider the inner amplitude.  On every component of $\{P<.2\}$,
\[
 (r^{11}A')'=r^{11}(P-V)A<0.
\]
At a positive left endpoint $A'\le0$, since the amplitude is entering
the sublevel set $A<.2^a$; at the origin $A'=0$.  Thus $A'<0$ at every
positive radius in this sublevel set.  If
$M=\max_{[0,r_b]}(A')_+$ is positive, it is attained at a point
$r_m\in(0,r_b)$: the derivative is zero at the origin and strictly
negative at $r_b$ for small $a$ by Lemma~\ref{profile:core-limit}.
At this point $P\ge.2$, $A''=0$, and $A'''\le0$.

Let $q=w(r_m)/r_m$.  The bounds $A\le.5^a$, $A(r_m)\ge.2^a$ and
\eqref{profile:w-ratio} give
\[
 .49<q\le q_+:=\frac{c_*}{12}(2.5)^{2a}=\frac12+O(a).
\]
The fixed lower bound follows from $h_b>.999$ and small $a$; it need
not approach $1/2$ at any prescribed rate.  Differentiating $V$ and
using \eqref{profile:amplitude} gives the exact identity
\[
 V'=r\left(\frac18-\frac{c_*}{2}\frac wr
                  +\frac{11}{2}\left(\frac wr\right)^2\right)
          +w^2\frac{A'}A.
\]
The polynomial $q\mapsto1/8-c_*q/2+11q^2/2$ is increasing for
$q\ge.49$ and small $a$, and its value at $q_+$ is $O(a)$, since its
value at $(a,q)=(0,1/2)$ is zero.  Thus at $r_m$,
\[
 V'\le Ca+w^2M/A.
\]
Differentiate the amplitude equation.  At this maximum,
\[
 (-\Delta A)'=-A'''-\frac{11}{r}A''+\frac{11}{r^2}A'\ge0,
\]
so
\[
 \frac{.2}{a}M\le VM+V'A\le CM+Ca.
\]
It follows that $M\le Ca^2$.  Since $w\ge0$ and $w/A$ is bounded on
the inner ball,
\[
 W(P/a)=\frac{wPA'}{a^2A}\le CP=Ca(P/a),
\]
which proves \eqref{profile:pressure-transport}.

In the exterior, \eqref{profile:tail} gives
$P/a=r^{-2}|f|^{1/a}/a$.  Repeated Euler differentiation and
$\kappa\le|f|\le\rho<1$ bound every fixed derivative by
$C_j a^{-N_j}\rho^{1/a}r^{-2}$ for some finite integer $N_j$.
This tends to zero uniformly after multiplication by $r^2$.
Also
$e^{2i\phi}=r^{4ib}f/\overline f$, so all its fixed Euler derivatives
are uniformly bounded.  As $w/r$ is uniformly bounded in the exterior,
these facts give \eqref{profile:exterior-pressure} and the exterior
bound on $W(P/a)$.  Together with $P\le.5$ in the interior they give
the uniform upper bound on its positive part.  By taking $a$ smaller
if necessary, the exterior formula also gives $P\le1$ everywhere.
The same bounds for $f$, and \eqref{profile:inner-bounds}, prove the
two-sided comparison for $\mu$ in \eqref{profile:coercive-bounds}.

It remains to prove positive deformation.  The two eigenvalues of the
derivative of a radial field are $w'$ and $w/r$.  For every limiting
matched profile they equal $1/2$ in the core.  On the exterior,
$A_0'<0$, $w_0>0$, and $w_0/r<.543$, so
\[
 w_0'=6-11w_0/r-2w_0A_0'/A_0\ge6-11(.543)=.027.
\]
Integrating from zero also gives $w_0/r\ge.027$.  The compact $C^1$
convergence of $A$ and formula \eqref{profile:w-ratio}, followed by
the last equation of \eqref{profile:amplitude}, give convergence of
both eigenvalues on every fixed compact interval, including zero.
For the remaining tail, the uniform expansion of Lemma~\ref{profile:exterior}
gives
\[
 \frac wr=\frac12+\frac{4b}{r^2}+O(r^{-4}),\qquad
 w'=\frac12-\frac{4b}{r^2}+O(r^{-4}),
\]
uniformly in small $a$.  Choose a large fixed radius to make these tail
bounds positive, and only then take $a$ small enough for the compact
convergence.  A subsequence contradiction, using compactness of the
parameter disk, makes this choice uniform for all matched profiles.
For example $c=.01$ is admissible after reducing $a_0$.  This proves
\eqref{profile:coercive-bounds}.
\end{proof}

\begin{proof}[Completion of Proposition~\ref{profile:existence}]
The inner construction, exterior construction and degree argument give
existence for every $0<a<a_0$ after finitely many reductions of $a_0$.
The order of choices is fixed: the disk and $\varepsilon$ first, then
the core comparison estimates, the exterior annulus and a single large
truncation order and starting radius, and finally sufficiently small
$a$.  Any later fixed derivative order is obtained by taking a higher
auxiliary truncation of the same exterior solution; it does not require
changing that solution or the already fixed geometric parameters.

The limiting statements on the inner ball are
Lemma~\ref{profile:core-limit}, and on $[r_b,\infty)$ follow from
Lemma~\ref{profile:exterior} and matching.  On a compact subset of
$(R,r_b)$, the limiting amplitude is strictly below $1$; uniform
$C^1$ convergence therefore places the profiles in a fixed annulus
bounded away from both $0$ and $1$.  All derivatives of
$z\mapsto|z|^{1/a}z$ tend to zero there.  The complex radial ODE,
successively differentiated, then proves $C^\infty$ convergence and
$P/a\to0$ on that compact set.  The weak-star pressure assertion on
any larger fixed ball follows by joining the inner assertion to this
exterior convergence.  The $C^1$ convergence of $\mu,w$ has already
been proved from their amplitude formulas.

Finally, radial Euler estimates imply Cartesian symbol estimates away
from zero: an ordinary derivative of order $j$ is a linear combination
of $r^{-j}$ times Euler derivatives with bounded smooth angular
coefficients.  The fixed-$a$ smoothness on a compact ball supplies the
remaining bounds in \eqref{profile:fixed-power-symbols}; their constants
are not asserted to be uniform as $a\downarrow0$.  As
$a-ib+\nu/2=0$,
\[
 (a-ib+D/2)Q=\tfrac12r^\nu Df,
\]
and \eqref{profile:symbol-expansion} gives
\eqref{profile:time-symbol}.  All parts of the proposition follow.
\end{proof}
